\documentclass[12pt,a4paper]{article}

% ------------------------------------------------------------
% Пакеты
% ------------------------------------------------------------
\usepackage[T2A]{fontenc}
\usepackage[utf8]{inputenc}
\usepackage[english,russian]{babel}

\usepackage{amsmath,amssymb,mathtools}
\usepackage{geometry}
\usepackage{graphicx}
\usepackage{booktabs}
\usepackage{array}
\usepackage{float}
\usepackage{enumitem}
\usepackage{hyperref}
\usepackage{xcolor}
\usepackage{fancyhdr}
\usepackage{titlesec}
\usepackage{caption}

\geometry{
    left=30mm,
    right=20mm,
    top=20mm,
    bottom=20mm
}

\hypersetup{
    colorlinks=true,
    linkcolor=black,
    citecolor=black,
    urlcolor=blue
}

\setlength{\parindent}{1.25cm}
\setlength{\parskip}{0pt}
\linespread{1.08}

\pagestyle{fancy}
\fancyhf{}
\fancyfoot[C]{\thepage}
\renewcommand{\headrulewidth}{0pt}

\titleformat{\section}
  {\normalfont\Large\bfseries}
  {\thesection.}{0.5em}{}

\titleformat{\subsection}
  {\normalfont\large\bfseries}
  {\thesubsection.}{0.5em}{}

\newcommand{\R}{\mathbb{R}}
\newcommand{\placeholder}[1]{\textbf{[PLACEHOLDER: #1]}}

\newcommand{\placeholderfigure}[2][0.8\textwidth]{
\begin{figure}[H]
    \centering
    \fbox{
        \parbox[c][5.5cm][c]{#1}{
            \centering
            \textbf{PLACEHOLDER}\\[0.5em]
            #2
        }
    }
\end{figure}
}

% ------------------------------------------------------------
% Документ
% ------------------------------------------------------------
\begin{document}

% ------------------------------------------------------------
% Титульный лист
% ------------------------------------------------------------
\begin{titlepage}
    \begin{center}
        {\small
        Федеральное государственное автономное образовательное учреждение\\
        высшего образования\\[0.3em]
        \textbf{«Национальный исследовательский университет ИТМО»}\\[0.8em]
        }

        \vfill

        {\Large \textbf{ОТЧЁТ}}\\[0.8em]
        {\large по проектной работе по дисциплине}\\[0.4em]
        {\large \textbf{«Математический анализ»}}\\[1.5em]

        {\LARGE
        \textbf{«Математический анализ в методах машинного обучения»}
        }

        \vfill

        \begin{flushright}
        \begin{tabular}{ll}
            Выполнили: &
            \begin{tabular}[t]{l}
                Вейгандт В.\\
                Зайцев Н.\\
                Горобец О.\\
                Козлов В.\\
                Шагеев Э.
            \end{tabular}
            \\[1.2em]
            Группа: & P 1.2\\[0.8em]
            Преподаватель: & Карпычев А.В.
        \end{tabular}
        \end{flushright}

        \vfill

        Санкт-Петербург\\
        2026
    \end{center}
\end{titlepage}

\tableofcontents
\newpage

% ------------------------------------------------------------
\section*{Введение}
\addcontentsline{toc}{section}{Введение}
% ------------------------------------------------------------

Машинное обучение рассматривает методы построения алгоритмов, поведение которых определяется не только явно заданными правилами, но и данными. В задачах обучения с учителем имеется набор объектов, для которых известны правильные ответы, и требуется построить модель, способную делать предсказания для новых объектов.

С математической точки зрения значительная часть таких задач естественно сводится к выбору функции из некоторого параметризованного семейства и поиску значений параметров, при которых ошибка модели минимальна. Поэтому при изучении машинного обучения возникают конструкции, непосредственно связанные с математическим анализом: функции нескольких переменных, частные производные, градиент, экстремумы, линии уровня и итерационные методы поиска минимума.

Цель настоящей работы --- рассмотреть основные идеи обучения с учителем и показать, как методы математического анализа возникают в них естественным образом. В качестве основных примеров рассматриваются линейная и логистическая регрессии, градиентные методы оптимизации, решающие деревья и градиентный бустинг.

Для достижения поставленной цели необходимо решить следующие задачи:
\begin{enumerate}[label=\arabic*)]
    \item сформулировать задачу обучения с учителем и ввести основные понятия машинного обучения;
    \item рассмотреть линейную регрессию как задачу минимизации среднеквадратичной ошибки;
    \item рассмотреть логистическую регрессию и вывести функцию потерь из принципа максимального правдоподобия;
    \item описать градиентный спуск и его основные модификации;
    \item рассмотреть принцип построения решающих деревьев и объяснить отличие этой задачи от гладкой параметрической оптимизации;
    \item рассмотреть градиентный бустинг как последовательное улучшение модели;
    \item реализовать основные рассматриваемые методы на Python с использованием NumPy и провести вычислительные эксперименты.
\end{enumerate}

Основой для теоретической части служат лекционные материалы по машинному обучению, учебные материалы Школы анализа данных Яндекса и исходный план проекта~\cite{shad-intro,shad-linreg,shad-logreg,project-plan,yandex-ml}. Для раздела, посвящённого современным градиентным оптимизаторам, дополнительно используются стандартные определения, согласующиеся с изложением в~\cite{d2l-opt}.

% ------------------------------------------------------------
\section{Постановка задачи машинного обучения}
% ------------------------------------------------------------

\subsection{Объекты, признаки и ответы}

Рассмотрим обучающую выборку
\[
    D = \{(x_i,y_i)\}_{i=1}^{n},
\]
где $x_i$ --- описание $i$-го объекта, а $y_i$ --- известный для него ответ.

Будем считать, что каждый объект представлен вектором численных признаков
\[
    x_i \in \R^d.
\]
Координаты этого вектора называются \emph{признаками} объекта. Множество допустимых значений ответа обозначим через $\mathcal{Y}$.

В зависимости от вида $\mathcal{Y}$ различают две основные задачи обучения с учителем.

\paragraph{Регрессия.}
Если
\[
    \mathcal{Y} = \R,
\]
требуется по признакам объекта предсказать вещественное число. Например, по характеристикам квартиры можно предсказывать её стоимость.

\paragraph{Классификация.}
Если
\[
    \mathcal{Y} = \{1,\ldots,K\},
\]
ответом является номер одного из $K$ классов. Например, банковская операция может относиться к классу обычных или подозрительных.

Таким образом, задача состоит в построении отображения
\[
    f:\R^d\to\mathcal{Y},
\]
которое хорошо восстанавливает зависимость между признаками объектов и ответами.

\subsection{Модель как семейство функций}

Искать решение среди всех возможных отображений практически невозможно. Поэтому заранее выбирается параметризованное семейство функций
\[
    \mathcal{F}
    =
    \{f_\theta:\theta\in\Theta\},
\]
где $\theta$ --- набор параметров модели.

Важно различать семейство моделей и конкретную обученную модель. Например, множество всех прямых на плоскости является семейством, а одна конкретная прямая соответствует конкретному набору параметров.

Обучение модели означает выбор такого значения параметров $\hat\theta$, которое определяется по обучающей выборке:
\[
    D
    \longmapsto
    \hat\theta
    \longmapsto
    f_{\hat\theta}.
\]

\subsection{Функция потерь}

Чтобы формально определить, какие параметры считать хорошими, вводится функция потерь
\[
    \ell:\mathcal{Y}\times\mathcal{Y}\to\R_{\ge 0},
\]
где $\ell(y,\hat y)$ измеряет ошибку одного предсказания.

Среднюю ошибку модели на обучающей выборке можно записать в виде
\[
    L(\theta)
    =
    \frac{1}{n}
    \sum_{i=1}^{n}
    \ell\bigl(y_i,f_\theta(x_i)\bigr).
\]

После фиксации данных $L$ зависит уже не от объектов выборки, а от параметров модели. Если
\[
    \theta=(\theta_1,\ldots,\theta_m)\in\R^m,
\]
то
\[
    L:\R^m\to\R
\]
является функцией нескольких переменных.

В этой постановке обучение параметрической модели принимает вид задачи оптимизации:
\[
    \boxed{
    \hat\theta\in\arg\min_{\theta\in\Theta}L(\theta)
    }.
\]

Тем самым основные понятия математического анализа возникают непосредственно из постановки задачи обучения: для поиска минимума необходимо исследовать поведение функции нескольких переменных.

\subsection{Обучающая и тестовая выборки}

Минимизация ошибки на тех же данных, на которых происходило обучение, сама по себе не гарантирует хорошей работы модели на новых объектах. Поэтому исходный набор данных обычно делят как минимум на две части:
\begin{itemize}
    \item \textbf{обучающую выборку} --- для подбора параметров модели;
    \item \textbf{тестовую выборку} --- для итоговой проверки качества на данных, которые не использовались при обучении.
\end{itemize}

При необходимости дополнительно выделяется валидационная выборка, на которой выбираются гиперпараметры и сравниваются варианты модели.

% ------------------------------------------------------------
\section{Линейные модели}
% ------------------------------------------------------------

\subsection{Линейная регрессия}

Одним из простейших параметрических семейств является семейство линейных функций. Для удобства свободный коэффициент можно включить в вектор параметров, добавив к вектору признаков единичную координату. Тогда
\[
    x\in\R^d,
    \qquad
    \theta\in\R^d,
\]
а предсказание имеет вид
\[
    \hat y
    =
    f_\theta(x)
    =
    x^T\theta.
\]

Для одного численного признака эта запись эквивалентна привычной прямой
\[
    \hat y = wx+b.
\]

При этом слово <<линейная>> относится к зависимости от параметров модели. При добавлении в признаки, например, величины $x^2$, модель
\[
    \hat y = \theta_0+\theta_1x+\theta_2x^2
\]
остаётся линейной по параметрам $\theta_0,\theta_1,\theta_2$.

\subsection{Среднеквадратичная ошибка}

Для регрессии естественно измерять отклонение предсказания от истинного ответа:
\[
    r_i
    =
    y_i-\hat y_i.
\]

Простое суммирование остатков неудобно, поскольку положительные и отрицательные ошибки могут взаимно компенсироваться. Один из стандартных вариантов --- использовать квадратичную функцию потерь
\[
    \ell(y,\hat y)
    =
    (y-\hat y)^2.
\]

Тогда среднеквадратичная ошибка имеет вид
\[
    L(\theta)
    =
    \frac{1}{n}
    \sum_{i=1}^{n}
    \bigl(y_i-x_i^T\theta\bigr)^2.
\]

Если составить матрицу признаков
\[
    X
    =
    \begin{pmatrix}
        x_1^T\\
        \vdots\\
        x_n^T
    \end{pmatrix}
    \in\R^{n\times d}
\]
и вектор ответов
\[
    Y
    =
    \begin{pmatrix}
        y_1\\
        \vdots\\
        y_n
    \end{pmatrix},
\]
то функцию потерь можно компактно записать:
\[
    L(\theta)
    =
    \frac{1}{n}
    \|Y-X\theta\|_2^2.
\]

В простейшем случае $\theta=(w,b)$ получаем функцию двух переменных
\[
    L=L(w,b).
\]
Её можно визуализировать как поверхность, а множества
\[
    L(w,b)=C
\]
--- как линии уровня. Минимум поверхности соответствует параметрам наилучшей прямой в смысле MSE.

\placeholderfigure{Поверхность $L(w,b)$ и линии уровня для линейной регрессии.}

\subsection{Градиент среднеквадратичной ошибки}

Для дальнейшего применения градиентных методов достаточно получить градиент MSE по параметрам модели. Для
\[
    L(\theta)
    =
    \frac{1}{n}
    \|Y-X\theta\|_2^2
\]
имеем
\[
    \boxed{
    \nabla_\theta L(\theta)
    =
    \frac{2}{n}
    X^T(X\theta-Y)
    }.
\]

Эта формула показывает, как изменение каждого параметра влияет на значение функции потерь. В частном случае линейной регрессии минимум MSE можно найти и аналитически, однако для общей логики работы далее важнее численный подход: имея значение градиента, можно последовательно изменять параметры в направлении уменьшения функции потерь.

\subsection{Логистическая регрессия}

Рассмотрим бинарную классификацию:
\[
    y_i\in\{0,1\}.
\]

Линейная комбинация признаков
\[
    s_\theta(x)=x^T\theta
\]
может принимать любые вещественные значения, поэтому напрямую интерпретировать её как вероятность нельзя.

В логистической регрессии вводится предположение о линейности логарифма шансов:
\[
    \log\frac{p_\theta(x)}{1-p_\theta(x)}
    =
    x^T\theta.
\]

Отсюда
\[
    \frac{p}{1-p}=e^{x^T\theta},
\]
поэтому
\[
    p_\theta(x)
    =
    \frac{1}{1+e^{-x^T\theta}}
    =
    \sigma(x^T\theta),
\]
где
\[
    \sigma(z)
    =
    \frac{1}{1+e^{-z}}
\]
--- сигмоидальная функция.

Таким образом,
\[
    p_\theta(x)
    =
    P_\theta(Y=1\mid x).
\]

При стандартном пороге $1/2$ класс определяется правилом
\[
    \hat y(x)
    =
    \begin{cases}
        1, & p_\theta(x)\ge \frac12,\\
        0, & p_\theta(x)<\frac12.
    \end{cases}
\]

\placeholderfigure{График сигмоиды и интерпретация вероятности класса.}

\subsection{Максимальное правдоподобие и cross-entropy}

Пусть
\[
    p_i
    =
    p_\theta(x_i).
\]
При условной независимости наблюдений вероятность получить конкретные ответы $y_1,\ldots,y_n$ при заданных параметрах $\theta$ записывается как
\[
    \mathcal{L}(\theta)
    =
    \prod_{i=1}^{n}
    p_i^{y_i}(1-p_i)^{1-y_i}.
\]

Метод максимального правдоподобия выбирает параметры
\[
    \hat\theta
    \in
    \arg\max_\theta
    \mathcal{L}(\theta).
\]

Поскольку логарифм является монотонно возрастающей функцией, максимум правдоподобия совпадает с максимумом логарифма правдоподобия:
\[
    \log\mathcal{L}(\theta)
    =
    \sum_{i=1}^{n}
    \left[
        y_i\log p_i
        +(1-y_i)\log(1-p_i)
    \right].
\]

Переходя от максимизации к минимизации отрицательного лог-правдоподобия, получаем функцию потерь
\[
    \boxed{
    L(\theta)
    =
    -\frac{1}{n}
    \sum_{i=1}^{n}
    \left[
        y_i\log p_i
        +(1-y_i)\log(1-p_i)
    \right]
    }.
\]

Это бинарная cross-entropy.

Несмотря на изменение модели и функции потерь, общая математическая постановка остаётся прежней:
\[
    \hat\theta
    \in
    \arg\min_\theta L(\theta).
\]

Именно поэтому один и тот же класс численных методов оптимизации может применяться к различным моделям.

% ------------------------------------------------------------
\section{Градиентные методы оптимизации}
% ------------------------------------------------------------

\subsection{Градиентный спуск}

Пусть
\[
    L:\R^d\to\R
\]
--- дифференцируемая функция потерь. Её градиент в точке $\theta$ определяется как
\[
    \nabla L(\theta)
    =
    \begin{pmatrix}
        \dfrac{\partial L}{\partial\theta_1}\\
        \vdots\\
        \dfrac{\partial L}{\partial\theta_d}
    \end{pmatrix}.
\]

Градиент задаёт направление наиболее быстрого локального роста функции. Поэтому для уменьшения значения $L$ естественно двигаться в противоположном направлении.

Итерация градиентного спуска:
\[
    \boxed{
    \theta_{t+1}
    =
    \theta_t
    -
    \eta_t\nabla L(\theta_t)
    },
\]
где $\eta_t>0$ называется \emph{скоростью обучения} или \emph{learning rate}.

На двумерной функции потерь траекторию алгоритма удобно изображать на линиях уровня.

\placeholderfigure{Траектория Gradient Descent на линиях уровня функции потерь.}

\subsection{Влияние размера шага}

Значение $\eta$ существенно влияет на поведение метода.

Рассмотрим одномерный пример
\[
    f(u)=u^2.
\]
Тогда
\[
    f'(u)=2u
\]
и итерация градиентного спуска имеет вид
\[
    u_{t+1}
    =
    u_t-2\eta u_t
    =
    (1-2\eta)u_t.
\]

Из этой формулы непосредственно видно, что при различных $\eta$ поведение последовательности может существенно отличаться: она может медленно сходиться, колебаться около минимума или расходиться.

В практической части рассматриваются три характерных режима работы градиентного спуска:
\begin{itemize}
    \item слишком маленький шаг;
    \item шаг, обеспечивающий устойчивую быструю сходимость;
    \item слишком большой шаг.
\end{itemize}

\placeholderfigure{Сравнение траекторий GD при различных значениях learning rate.}

\subsection{Momentum}

На вытянутых поверхностях уровня обычный градиентный спуск может совершать сильные колебания поперёк узкой долины. Метод Momentum добавляет к текущему направлению информацию о предыдущих шагах:
\[
    v_t
    =
    \beta v_{t-1}
    +
    g_t,
    \qquad
    g_t=\nabla L(\theta_t),
\]
\[
    \theta_{t+1}
    =
    \theta_t
    -
    \eta v_t.
\]

Параметр $\beta\in[0,1)$ определяет степень сохранения накопленного направления. За счёт этого последовательные градиенты усиливают движение вдоль устойчивого направления и частично компенсируют поперечные колебания.

\subsection{Nesterov Momentum}

Метод Нестерова использует идею ``заглядывания вперёд'': градиент оценивается не только в текущей точке, но с учётом предполагаемого движения по накопленному импульсу.

Одна из стандартных форм записи:
\[
    g_t
    =
    \nabla L
    \left(
        \theta_t-\eta\beta v_{t-1}
    \right),
\]
\[
    v_t
    =
    \beta v_{t-1}
    +
    g_t,
\]
\[
    \theta_{t+1}
    =
    \theta_t-\eta v_t.
\]

В результате алгоритм раньше реагирует на изменение направления поверхности функции потерь.

\subsection{AdaGrad}

В многомерных задачах функция может изменяться с различной скоростью по разным координатам. AdaGrad использует отдельный эффективный размер шага для каждой координаты.

Пусть
\[
    g_t
    =
    \nabla L(\theta_t).
\]
Накопим суммы квадратов градиентов:
\[
    G_t
    =
    G_{t-1}
    +
    g_t\odot g_t,
\]
где $\odot$ обозначает покоординатное умножение.

Обновление параметров:
\[
    \theta_{t+1}
    =
    \theta_t
    -
    \eta
    \frac{g_t}
         {\sqrt{G_t}+\varepsilon},
\]
где деление и квадратный корень выполняются покоординатно.

Недостаток метода состоит в том, что $G_t$ постоянно возрастает, поэтому эффективный шаг со временем может стать слишком малым.

\subsection{RMSProp}

RMSProp заменяет полное накопление квадратов градиентов экспоненциальным скользящим средним:
\[
    s_t
    =
    \rho s_{t-1}
    +
    (1-\rho)g_t\odot g_t.
\]

Параметры обновляются по правилу
\[
    \theta_{t+1}
    =
    \theta_t
    -
    \eta
    \frac{g_t}
         {\sqrt{s_t}+\varepsilon}.
\]

В отличие от AdaGrad, старые значения постепенно теряют влияние, поэтому шаг не обязан монотонно стремиться к нулю.

\subsection{Adam}

Adam объединяет идеи Momentum и RMSProp. Вводятся экспоненциальные средние первого и второго моментов градиента:
\[
    m_t
    =
    \beta_1m_{t-1}
    +
    (1-\beta_1)g_t,
\]
\[
    v_t
    =
    \beta_2v_{t-1}
    +
    (1-\beta_2)g_t\odot g_t.
\]

Так как в начале $m_0=v_0=0$, применяются поправки на смещение:
\[
    \hat m_t
    =
    \frac{m_t}{1-\beta_1^t},
    \qquad
    \hat v_t
    =
    \frac{v_t}{1-\beta_2^t}.
\]

Итоговое обновление:
\[
    \boxed{
    \theta_{t+1}
    =
    \theta_t
    -
    \eta
    \frac{\hat m_t}
         {\sqrt{\hat v_t}+\varepsilon}
    }.
\]

Таким образом, Adam одновременно использует накопленное направление движения и адаптивный масштаб шага по отдельным координатам.

\subsection{Сравнение методов}

Для корректного визуального сравнения рассматриваемые оптимизаторы запускаются:
\begin{itemize}
    \item на одной функции;
    \item из одной стартовой точки;
    \item при заранее зафиксированных гиперпараметрах.
\end{itemize}

Сравниваются следующие методы:
\[
    \text{GD},
    \quad
    \text{Momentum},
    \quad
    \text{Nesterov},
    \quad
    \text{AdaGrad},
    \quad
    \text{RMSProp},
    \quad
    \text{Adam}.
\]

Основными объектами сравнения являются траектории на линиях уровня и изменение функции потерь по итерациям.

\placeholderfigure{Траектории шести оптимизаторов из одной стартовой точки.}

% ------------------------------------------------------------
\section{Решающие деревья и градиентный бустинг}
% ------------------------------------------------------------

\subsection{Решающие деревья}

Решающее дерево строит предсказание последовательностью простых условий. Для вещественного признака типичный предикат имеет вид
\[
    x_j \le t,
\]
где $j$ --- номер признака, а $t$ --- порог.

Пусть в некоторую вершину дерева пришло множество объектов $S$. Кандидат на разбиение $(j,t)$ делит его на две части:
\[
    S_L(j,t)
    =
    \{(x_i,y_i)\in S:\ x_{ij}\le t\},
\]
\[
    S_R(j,t)
    =
    \{(x_i,y_i)\in S:\ x_{ij}>t\}.
\]

Задача алгоритма в этой вершине --- выбрать такое разбиение, после которого объекты в дочерних вершинах становятся как можно более однородными по ответам.

\subsection{Критерий разбиения}

Обозначим через $H(S)$ меру неоднородности множества объектов в вершине. Для разбиения $(j,t)$ можно определить качество как уменьшение неоднородности:
\[
    Q(S,j,t)
    =
    H(S)
    -
    \frac{|S_L|}{|S|}H(S_L)
    -
    \frac{|S_R|}{|S|}H(S_R).
\]

На каждом шаге дерево выбирает
\[
    (j^\ast,t^\ast)
    \in
    \arg\max_{j,t} Q(S,j,t).
\]

Таким образом, рассматриваются возможные признаки и пороги, после чего жадно выбирается локально лучшее разбиение.

\paragraph{Регрессия.}
В задаче регрессии в качестве неоднородности естественно использовать среднеквадратичное отклонение ответов от их среднего:
\[
    H_{\mathrm{MSE}}(S)
    =
    \frac{1}{|S|}
    \sum_{(x_i,y_i)\in S}
    (y_i-\bar y_S)^2,
\]
где
\[
    \bar y_S
    =
    \frac{1}{|S|}
    \sum_{(x_i,y_i)\in S}y_i.
\]

Эквивалентно можно выбирать признак $j$ и порог $t$, минимизирующие сумму квадратов ошибок в двух дочерних вершинах:
\[
    \min_{j,t}
    \left[
        \sum_{(x_i,y_i)\in S_L}
        (y_i-\bar y_L)^2
        +
        \sum_{(x_i,y_i)\in S_R}
        (y_i-\bar y_R)^2
    \right].
\]

Предсказанием в листе регрессионного дерева обычно служит среднее значение ответов объектов, попавших в этот лист.

\paragraph{Классификация.}
Для классификации часто используют критерий Джини
\[
    G(S)
    =
    1-\sum_{k=1}^{K}p_k^2
\]
или энтропию
\[
    H(S)
    =
    -\sum_{k=1}^{K}p_k\log p_k,
\]
где $p_k$ --- доля объектов класса $k$ в вершине.

Чем меньше неоднородность дочерних вершин, тем удачнее разбиение. В листе классификационного дерева можно возвращать наиболее частый класс или оценки вероятностей классов по их долям среди объектов листа.

\subsection{Жадное построение дерева}

Полный перебор всех возможных структур дерева практически невозможен, поэтому дерево строится жадно:
\begin{enumerate}
    \item в текущую вершину приходит множество объектов $S$;
    \item перебираются допустимые признаки $j$ и пороги $t$;
    \item для каждого кандидата вычисляется значение критерия разбиения;
    \item выбирается лучшее разбиение $(j^\ast,t^\ast)$;
    \item процедура рекурсивно повторяется для левой и правой дочерних вершин.
\end{enumerate}

Рост дерева останавливают, например, при достижении максимальной глубины, слишком малом числе объектов в вершине или недостаточном уменьшении неоднородности. Эти ограничения позволяют контролировать сложность дерева и уменьшать переобучение.

\subsection{Почему обычный Gradient Descent здесь неприменим}

В линейных моделях параметры принадлежат непрерывному пространству $\R^d$, а функция потерь рассматривается как функция этих параметров.

При построении дерева необходимо выбирать номер признака, порог и саму дискретную структуру дерева. Поэтому нельзя непосредственно вычислить производную функции потерь, например, по номеру выбранного признака. Вместо движения по антиградиенту используется жадный перебор допустимых локальных разбиений.

\subsection{Случайность и разнообразие деревьев}

Для \emph{одного} обычного решающего дерева случайный выбор признаков не является обязательной частью алгоритма: при поиске split'а можно рассматривать все доступные признаки.

Случайный выбор объектов и признаков становится особенно важен в ансамблях независимых деревьев. Например, в Random Forest каждое дерево строится на собственной bootstrap-выборке объектов, а при поиске разбиения в вершине рассматривается случайное подмножество признаков. Благодаря этому отдельные деревья получаются менее похожими друг на друга, а усреднение их предсказаний уменьшает разброс ансамбля.

В Gradient Boosting деревья различаются по другой причине: каждое следующее дерево обучается уже на новой задаче, зависящей от ошибок текущего ансамбля. Дополнительный случайный subsampling объектов или признаков также может применяться, но не является обязательным для самой идеи градиентного бустинга.

\subsection{Идея ансамбля}

Вместо одной модели можно построить композицию нескольких базовых алгоритмов:
\[
    F(x)
    =
    \sum_{m=1}^{M}\alpha_m h_m(x),
\]
где $h_m$ --- отдельные базовые модели.

В Random Forest деревья строятся в значительной степени независимо, после чего их ответы усредняются или объединяются голосованием. Градиентный бустинг использует другую схему: деревья строятся последовательно, и каждый следующий алгоритм старается исправить ошибку уже построенной композиции.

\subsection{Gradient Boosting}

Пусть после $m-1$ шагов имеется модель
\[
    F_{m-1}(x).
\]
На следующем шаге к ней добавляется новое дерево:
\[
    F_m(x)
    =
    F_{m-1}(x)
    +
    \eta h_m(x),
\]
где $\eta$ --- размер шага бустинга.

Для квадратичной функции потерь
\[
    \ell(y,F)
    =
    (y-F)^2
\]
отрицательный градиент по текущему предсказанию пропорционален обычному остатку
\[
    y_i-F_{m-1}(x_i).
\]
Поэтому в этом случае очередное дерево можно обучать приближать residuals текущего ансамбля.

В общем случае вводятся псевдоостатки
\[
    r_i^{(m)}
    =
    -
    \left.
    \frac{\partial
        \ell(y_i,F(x_i))}
        {\partial F(x_i)}
    \right|_{F=F_{m-1}}.
\]

Новое дерево $h_m$ обучается приближать эти значения, после чего добавляется к ансамблю. Тем самым идея антиградиента переносится с пространства обычных параметров на последовательное построение самой функции $F$.

\placeholderfigure{Пошаговое добавление деревьев и изменение residuals.}

% ------------------------------------------------------------
\section{Дополнительные методы}
% ------------------------------------------------------------

Основная часть работы посвящена обучению с учителем и градиентным методам. Дополнительно рассматриваются несколько подходов, использующих другие математические идеи~\cite{project-plan}.

\subsection{Principal Component Analysis}

PCA применяется для уменьшения размерности данных. Пусть строки матрицы $X$ содержат центрированные объекты. Рассмотрим ковариационную матрицу
\[
    C
    =
    \frac{1}{N}X^TX.
\]

Первая главная компонента определяется направлением единичной длины, вдоль которого дисперсия проекции максимальна:
\[
    \max_{\|w\|=1} w^TCw.
\]

Решением является собственный вектор матрицы $C$, соответствующий наибольшему собственному значению. Геометрически PCA выбирает направления, вдоль которых данные обладают наибольшим разбросом, и позволяет отбросить направления с относительно небольшой вариативностью.

\subsection{Метод Монте-Карло}

Методы Монте-Карло используют случайную выборку для приближённого вычисления математических величин. Например,
\[
    I
    =
    \int_a^b f(x)\,dx.
\]

Если
\[
    X\sim U(a,b),
\]
то
\[
    I
    =
    (b-a)\mathbb{E}[f(X)].
\]

Поэтому интеграл можно приближённо оценить по случайным точкам:
\[
    I
    \approx
    \frac{b-a}{N}
    \sum_{i=1}^{N}f(x_i).
\]

Эта идея важна не только для численного интегрирования: случайные оценки математических ожиданий постоянно возникают в статистике, оптимизации и машинном обучении.

\subsection{Муравьиный алгоритм}

Ant Colony Optimization является вероятностным метаэвристическим методом оптимизации. Он не относится непосредственно к рассматриваемым выше алгоритмам обучения с учителем, но служит примером другого подхода к поиску сложного решения.

Вероятность перехода из вершины $i$ в вершину $j$ можно задавать в виде
\[
    P_{ij}
    =
    \frac{\tau_{ij}^{\alpha}\eta_{ij}^{\beta}}
    {\sum\limits_{k\in A_i}\tau_{ik}^{\alpha}\eta_{ik}^{\beta}},
\]
где $\tau_{ij}$ описывает накопленный феромон, $\eta_{ij}$ --- локальную привлекательность перехода, а $A_i$ --- множество допустимых переходов.

После итерации феромон обновляется:
\[
    \tau_{ij}
    \leftarrow
    (1-\rho)\tau_{ij}
    +
    \Delta\tau_{ij}.
\]

Работа метода демонстрируется на небольшой задаче коммивояжёра, где отслеживаются изменение феромонов, вероятностей переходов и найденного маршрута.

% ------------------------------------------------------------
\section{Практическая часть}
% ------------------------------------------------------------


% ------------------------------------------------------------

\subsection{Используемые средства}

Практическая часть реализуется на языке Python. Основными используемыми библиотеками являются:
\begin{itemize}
    \item \texttt{NumPy} --- для вычислений с массивами и реализации алгоритмов;
    \item \texttt{Matplotlib} --- для построения графиков и визуализации траекторий;
    \item \texttt{Pandas} --- при необходимости для удобного представления табличных данных.
\end{itemize}

Основные алгоритмы обучения реализуются самостоятельно. Готовые реализации из \texttt{scikit-learn} используются только для дополнительной проверки результатов и не заменяют собственную реализацию исследуемых алгоритмов.

\subsection{Эксперимент 1: поверхность loss линейной регрессии}

Для простого набора точек рассматривается модель
\[
    \hat y=wx+b
\]
и функция
\[
    L(w,b)
    =
    \frac1n
    \sum_{i=1}^{n}
    \bigl(y_i-wx_i-b\bigr)^2.
\]

В ходе эксперимента выполняются следующие действия:
\begin{enumerate}
    \item построить исходные данные;
    \item построить поверхность $L(w,b)$;
    \item показать линии уровня;
    \item отметить положение минимума;
    \item визуализировать изменение прямой при движении параметров к минимуму.
\end{enumerate}

\placeholderfigure{Итоговая поверхность MSE и положение найденного минимума.}

\paragraph{Результат эксперимента.}
\placeholder{результаты эксперимента и численные значения}

\subsection{Эксперимент 2: влияние learning rate}

Из одной стартовой точки запускается Gradient Descent с несколькими значениями $\eta$.

Рассматриваются три типичных режима:
\begin{itemize}
    \item медленная сходимость при малом шаге;
    \item устойчивая сходимость при подходящем шаге;
    \item колебания или расходимость при слишком большом шаге.
\end{itemize}

\placeholderfigure{Сравнение трёх траекторий Gradient Descent.}

\paragraph{Результат эксперимента.}
\placeholder{значения $\eta$ и анализ полученных траекторий}

\subsection{Эксперимент 3: сравнение оптимизаторов}

На одной двумерной функции из одной начальной точки запускаются:
\[
    \text{GD},
    \text{Momentum},
    \text{Nesterov},
    \text{AdaGrad},
    \text{RMSProp},
    \text{Adam}.
\]

Для каждого метода фиксируются:
\begin{itemize}
    \item число итераций;
    \item траектория;
    \item значения функции потерь;
    \item расстояние до известного минимума, если оно может быть вычислено.
\end{itemize}

\placeholderfigure{Общий график траекторий оптимизаторов.}

\begin{table}[H]
    \centering
    \caption{Сравнение оптимизаторов}
    \begin{tabular}{lccc}
        \toprule
        Метод & Итерации & Финальный loss & Комментарий \\
        \midrule
        GD       & PLACEHOLDER & PLACEHOLDER & PLACEHOLDER \\
        Momentum & PLACEHOLDER & PLACEHOLDER & PLACEHOLDER \\
        Nesterov & PLACEHOLDER & PLACEHOLDER & PLACEHOLDER \\
        AdaGrad  & PLACEHOLDER & PLACEHOLDER & PLACEHOLDER \\
        RMSProp  & PLACEHOLDER & PLACEHOLDER & PLACEHOLDER \\
        Adam     & PLACEHOLDER & PLACEHOLDER & PLACEHOLDER \\
        \bottomrule
    \end{tabular}
\end{table}

\subsection{Эксперимент 4: построение небольшого дерева}

На небольшом наборе данных воспроизводятся несколько первых шагов построения дерева:
\begin{enumerate}
    \item перебрать возможные признаки и пороги;
    \item вычислить значение критерия для каждого кандидата;
    \item выбрать лучшее разбиение;
    \item повторить процедуру в дочерних вершинах.
\end{enumerate}

\placeholderfigure{Исходные данные и выбранные границы разбиения дерева.}

\paragraph{Результат эксперимента.}
\placeholder{выбранные разбиения и значения критерия}

\subsection{Эксперимент 5: пошаговый Gradient Boosting}

Для демонстрации работы градиентного бустинга рассматривается последовательность
\[
    F_0,
    F_1,
    \ldots,
    F_M.
\]

На каждом шаге сохраняются:
\begin{itemize}
    \item текущие предсказания;
    \item residuals или pseudo-residuals;
    \item новое дерево;
    \item значение общей функции потерь.
\end{itemize}

Итоговая визуализация показывает изменение функции потерь при увеличении числа деревьев.

\placeholderfigure{График loss в зависимости от числа добавленных деревьев.}

\paragraph{Результат эксперимента.}
\placeholder{результаты эксперимента}

% ------------------------------------------------------------

\subsection{Дополнительные демонстрации}

В качестве дополнительных демонстраций рассматриваются следующие визуализации:
\begin{itemize}
    \item PCA: двумерное облако точек, главные направления и проекция на первую компоненту;
    \item Monte Carlo: приближение интеграла при возрастающем числе случайных точек;
    \item Ant Colony Optimization: изменение феромонов и маршрута в небольшой задаче коммивояжёра.
\end{itemize}

Дополнительные демонстрации иллюстрируют методы, основанные на других математических идеях, и дополняют основную практическую часть по линейным моделям, оптимизаторам, деревьям и градиентному бустингу.

% ------------------------------------------------------------
\section{Распределение работы между участниками}
% ------------------------------------------------------------

Работа над проектом распределена между участниками следующим образом:

\begin{itemize}
    \item \textbf{Зайцев Н.} --- составление отчёта и подготовка презентации;
    \item \textbf{Вейгандт В.} --- написание программного кода и реализация рассматриваемых алгоритмов;
    \item \textbf{Козлов В.} --- создание анимаций и визуализаций, сопровождающих основные алгоритмы;
    \item \textbf{Горобец О.} --- подготовка и вывод математических доказательств и выкладок;
    \item \textbf{Шагеев Э.} --- танец продматника.
\end{itemize}

% ------------------------------------------------------------
\section*{Заключение}
\addcontentsline{toc}{section}{Заключение}
% ------------------------------------------------------------

В работе была сформулирована базовая задача обучения с учителем и рассмотрены несколько ключевых методов классического машинного обучения.

На примере линейной регрессии было показано, что обучение параметрической модели может быть представлено как поиск минимума функции нескольких переменных. Среднеквадратичная ошибка задаёт функцию от параметров модели, а её градиент определяет направление локального изменения ошибки и естественно приводит к численным методам оптимизации.

Для задачи бинарной классификации была рассмотрена логистическая регрессия. Функция cross-entropy была получена из принципа максимального правдоподобия, после чего задача обучения вновь свелась к минимизации функции потерь.

Градиентный спуск был рассмотрен как общий численный метод минимизации. Его модификации --- Momentum, Nesterov, AdaGrad, RMSProp и Adam --- изменяют способ использования информации о градиентах и позволяют улучшать движение по сложным поверхностям функции потерь.

Решающие деревья представляют другой тип моделей: их структура выбирается в дискретном пространстве, поэтому обычный градиентный спуск непосредственно к построению дерева не применяется. Был подробно рассмотрен жадный выбор разбиений и роль критериев неоднородности. Также было отмечено, что случайный выбор объектов и признаков характерен прежде всего для Random Forest, тогда как в Gradient Boosting разнообразие деревьев возникает за счёт последовательного исправления ошибок текущего ансамбля.

Дополнительно были кратко рассмотрены PCA, метод Монте-Карло и муравьиный алгоритм как примеры применения линейной алгебры, вероятностных идей и метаэвристической оптимизации.

Таким образом, математический анализ не является отдельным внешним дополнением к рассматриваемым алгоритмам. Понятия функции нескольких переменных, частных производных, градиента, экстремума и сходимости естественно возникают непосредственно из задачи обучения моделей.

Вычислительная часть дополняет теоретическое рассмотрение экспериментальным сравнением поведения изученных методов.

% ------------------------------------------------------------
\begin{thebibliography}{99}
\addcontentsline{toc}{section}{Список источников}

\bibitem{shad-intro}
Броиловский А.
\textit{Лекция 1. Обзор задач машинного обучения}.
По мотивам лекции Н. Волкова, ШАД СПб, 2026.

\bibitem{shad-linreg}
Овчаров С.
\textit{Лекция 2. Линейная регрессия и основы оптимизации}.
По мотивам лекции Н. Волкова, ШАД СПб, 2026.

\bibitem{shad-logreg}
Овчаров С.
\textit{Лекция 3. Логистическая регрессия и классификация}.
По мотивам лекции Н. Волкова, ШАД СПб, 2026.

\bibitem{project-plan}
\textit{Математический анализ в методах машинного обучения}.
Материалы проекта, 2026.

\bibitem{yandex-ml}
Школа анализа данных Яндекса.
\textit{Учебник по машинному обучению}.
\url{https://contest.yandex.ru/tracks/ml/}.

\bibitem{yandex-linear}
Синицин Ф., Соколов Е.
\textit{Линейные модели}.
Учебник по машинному обучению ШАД Яндекса.
\url{https://contest.yandex.ru/tracks/ml/supervised-learning/linear-models}.

\bibitem{d2l-opt}
Zhang A., Lipton Z. C., Li M., Smola A. J.
\textit{Dive into Deep Learning. Chapter 12: Optimization Algorithms}.
\url{https://d2l.ai/chapter_optimization/index.html}.


\bibitem{yandex-trees}
Школа анализа данных Яндекса.
\textit{Решающие деревья}.
\url{https://contest.yandex.ru/tracks/ml/supervised-learning/decision-trees}.

\bibitem{yandex-boosting}
Школа анализа данных Яндекса.
\textit{Градиентный бустинг}.
\url{https://contest.yandex.ru/tracks/ml/supervised-learning/gradient-boosting}.

\bibitem{sklearn-ensembles}
Scikit-learn Developers.
\textit{Ensembles: Gradient boosting, random forests, bagging, voting, stacking}.
\url{https://scikit-learn.org/stable/modules/ensemble.html}.

\end{thebibliography}

% ------------------------------------------------------------
\appendix
\section{Программная реализация}
% ------------------------------------------------------------

Исходный код проекта доступен в репозитории:

\begin{center}
    \placeholder{ссылка на репозиторий}
\end{center}

\end{document}
